\glsresetall

\chapter{Enhaced Physical Memory Protection}
% \Gls{pmp} is relatively permissive and more restriction are wanted.
% To add more restrictions \gls{smepmp} was developed.
% This adds the possibility for Machine-mode only memory with access restrictions as well as User- and Supervisor-mode only memory.
% It also adds the possibility to disallow memory access when memory is not within a \gls{pmp} region, including in Machine-mode.
% While \gls{pmp} effectively isolates User- and Supervisormode from Machine-mode, it assumes Machine-mode code is trustworthy.
% \Gls{smepmp} addresses this by adding restriction for potentially buggy Machine-mode code~\cite{smePMP_apendix}.
\Gls{pmp} is relatively permissive and assumes that Machine-mode software is fully trusted.
However, additional restrictions on Machine-mode memory accesses are desirable to protect against potentially buggy or vulnerable Machine-mode code.
To address this, the RISC-V specification introduces \gls{smepmp}.

\Gls{smepmp} extends \gls{pmp} with additional protection mechanisms.
It introduces the possibility to define Machine-mode-only memory regions with restricted access permissions, as well as memory regions accessible only from User- and Supervisor-mode.
Furthermore, it allows memory accesses outside configured \gls{pmp} regions to be denied even in Machine-mode.
While \gls{pmp} isolates User- and Supervisor-mode from Machine-mode, \gls{smepmp} additionally enables restrictions within Machine-mode itself.~\cite{smePMP_apendix}

\section{\texttt{mseccfg} Control Status Register}
To configure these additional security mechanisms, \gls{smepmp} introduces several control bits in the \texttt{mseccfg} \gls{csr}.
The relevant fields are \gls{mmwp}, \gls{mml}, and \gls{rlb}.

The existing Katamaran \gls{pmp} model does not include the \texttt{mseccfg} register, since the currently implemented features do not require it.
Supporting \gls{smepmp} therefore requires extending the model with the \texttt{mseccfg} register together with definitions for accessing its individual fields by name.

% This \gls{csr} is added by adding it to the list of existing \glspl{csr}.
% We then need to define the layout of this \gls{csr} which also defines the names of the seperate fields.
% Now that the register and its fields are defined we need to ensure the proofs also know these exist to verify the correct working.
% In total this is a change of 75 lines in 5 files.

Support for this \gls{csr} is added by adding the \texttt{mseccfg} register to the list of existing \glspl{csr}.
The register layout and associated field definition were then added to the model, this allows for the use of the bit by name.
Finally, the proof's helper functions are updated to accommodate the new register and its fields.
% explain that the proof (specification) need to know onwership of the register

Adding support for the \texttt{mseccfg} register requires modifications across five files and introduce approximatly 75 lines of changes.
% FIXME move amount changed to evaluation
% TODO how much changed

\section{Changes to PMP Semantics}
The fields in the \texttt{mseccfg} register extend the behavior of the existing \gls{pmp} mechanism.
As a result, the access-checking semantics implemented in the current Katamaran \gls{pmp} model must be updated to correctly represent the \gls{smepmp} specification.

The following subsections describe the behavior introduced by each \texttt{mseccfg} field together with the corresponding changes required in the formal model.

\subsection{Machine-Mode Whitelist Policy}
The \gls{mmwp} bit is the bit that disables the possibility for Machine-mode to access memory that is not covered by any \gls{pmp} range.
This ensures that all accessible memory regions must be explicitly configured, forcing the system designer to intentionally specify the required memory permissions.
The \gls{mmwp} bit is also a sticky bit, meaning that once this bit is set it cannot be reset until the whole system resets.
This prevents the restriction from being circumvented by temporarily disabling the bit.

The Katamaran \gls{pmp} model determines which \gls{pmp} region applies to a memory access.
With \gls{pmp}, accesses from Machine-mode are permitted when the accessed memory does not fall within any \gls{pmp} range while access from other modes are denied.
When \gls{mmwp} is set, this behavior chages such that Machine-mode can no longer access unmatched memory.

\subsection{Machine-Mode Lockdown}
\Gls{mml} is the bit that allows for the creation of restricted machine only memory and the creation of User- and Supervisor-mode only memory.
When this bit is set it adds functionallity to the L-bit in the \texttt{pmpcfg} registers.
With the \gls{mml} bit set a locked entry becomes Machine-mode only, while an entry that is not locked becomes User- and Supervisor mode only.
This bit, just like \gls{mmwp} is a sticky bit to prevent circumvention by disabling the bit.

This does however not allow for shared memory between Machine-mode and User- and Supervisor-mode.
To solve this the previously reserved combinations with the read-bit being `0' and the write bit `1' are used to be shared memory.
The combination of both the lock-, read-, write- and execute- bit being `1' also becomes shared memory.
The permissions of these combinations also do not reflect the read-, write- and execute- bits but follow the permissions in~\cref{fig:smepmp_visual}. % TODO move image to this chapter

%TODO changes in katamaran

\subsection{Rule Locking Bypass}
\Gls{rlb} is a bypass of the locking rule made to fix an issue caused by \gls{mml} as described in the RISC-V \gls{smepmp} appendix:
``There are use cases where developers want to enforce PMP rules in M-mode during the boot process, that are also able to modify, merge, and / or remove later on.
Since a rule that is enforced in M-mode also needs to be locked (or else badly written or malicious M-mode software can remove it at any time),
the only way for developers to approach this is to keep adding PMP rules to the chain and rely on rule priority.
This is a waste of PMP rules and since it's only needed during boot, mseccfg.RLB is a simple workaround that can be used temporarily and then disabled and locked down''.~\cite{smePMP_apendix}
The \gls{rlb} rule thus makes it possible to modify locked \gls{pmp} entries.

To ensure this is not used to bypass entries that do need to be locked, this bit can only be set when no entries are locked.

%TODO changes in katamaran

\section{Changes in the Proofs}

\section{some conclusion}% TODO better title
